\scanpage{308}
\mspage{224}

\margin{Ce lemme fait partie des \unsure{réslts|résultats} généraux sur la cofinalité et les cat. filtrantes.}% slanted

\uline{\deleted{\unsure{Proposition}}} 4.9.6.3. Soit $f\colon I^{\unsure{a|4}}\unsure{\text{---}|\longrightarrow}I'$ un foncteur, et considérons les conditions :

a) $I$ filtrante, $f$ cofinal i.e. \inserted{(en présence de $I$ \unsure{filtr.te|filtr.-e})} $f$ satisfait les conditions fonctorielles $\text{\textcircled{F1}},\text{\textcircled{F2}}$.

\inserted{b) $I'$ est filtrante, et les $\vphantom{I}_{i'}\backslash I$ (pour $i'\in\operatorname{Ob}I'$) sont filtrantes.}

\unsure{c)|b)} $I'$ filtrante, et $f$ satisfait à F1, ainsi qu'à la condition suivante

\textcircled{F3} Pour tout couple $(i_0,i_1)$ d'objets de $I$ et toute flèche \deleted{\unsure{$u$:}} $f(i_0)\xrightarrow{u'}f(i_1)$ il existe un objet $i_2$ de $I$, et une flèche $v'\colon f(i_1)\longrightarrow f(i_2)$, tels que \deleted{\unsure{les|des}} $v'$ et $v'u'$ \deleted{\unsure{puissent s'écrire}} \deleted{\unsure{proviennent des}} \inserted{soient relevables} \unsure{en|des} flèches \deleted{\unsure{de}} $i_{\unsure{1|0}}\xrightarrow{v}i_2$ et $i_0\overset{w}{\unsure{\text{---}|\longrightarrow}}i_2$ \unsure{;|:}

\diagram{
\begin{tikzcd}[ampersand replacement=\&]
\unsure{i_0|j_0} \arrow[r,"{\deleted{\unsure{u}}}",dotted] \arrow[rr,"w",bend left=21,dotted] \arrow[d,no head] \& i_1 \arrow[r,"v",dotted] \& i_2 \\
f(i_0) \arrow[r,"{u'}"] \arrow[rr,"{f(w)}"',bend right=22,dashed] \& f(i_1) \arrow[r,"{f(v)=v'}",no head] \& f(i_2)
\end{tikzcd}
}

$[$ i.e il existe $\unsure{i_{2a}|i_2}$ et $v\colon i_1\unsure{\longrightarrow|\text{---}}i_2$, $w\colon i_0\unsure{\longrightarrow|\text{---}}i_2$ tels qu'on ait $f(v)u'=f(w)$ $]$.

\deleted{\unsure{c')} $I'$ filtrante, les \unsure{$I_{i'}$} et les $\vphantom{I}_{i'}\backslash I$ \unsure{ps.|pseudo-}filt.}

\inserted{b') $I'$ $0$-connexe, les $\vphantom{I}_{i'}\backslash I$ filtrantes.}

\inserted{b'') $I'$ est filtrante, \inserted{satisfait PF$(2)^*$ (égalisation des doubles flèches)} les $\vphantom{I}_{i'}\backslash I$ ps. filtr\unsure{antes|.}, \deleted{\unsure{les}} \unsure{i.e.} satisfont à la condition PF$1'$.}

\margin{p. ex. PF\unsure{3|2} est satisfaite si $I'$ est \uline{ordonnée}, sans avoir à supposer alors de $I'$ filtrante, sauf \deleted{\unsure{ment}} \uline{$0$-connexe}}% slanted

Ceci posé :

1) On a les implications

\[
\begin{aligned}&\deleted{\unsure{a\Longleftrightarrow b|a\Longleftarrow b}}\unsure{\Longrightarrow|\Longleftrightarrow}c\\
&a\Longleftrightarrow b\Longleftrightarrow b'\Longrightarrow b''\unsure{\Longleftrightarrow|\Longleftarrow}c\quad\deleted{\unsure{b\Longleftrightarrow c|\Longleftrightarrow a}}\end{aligned}
\]

et si $f$ est \uline{fidèle} \inserted{(pour $I$ \unsure{ordonnée|ndonnée})}, les \deleted{trois} \inserted{\unsure{cinq|quatre}} conditions a), b), \inserted{b'), b''),} c) sont équivalentes.

\deleted{2) \unsure{plus|pour} \unsure{$\pi$|R} \unsure{un|un card.} \unsure{cardinal infini :} \unsure{et} condition : \unsure{équiv|équivalente|Supposons}} \deleted{\unsure{Pour que} Supposons $I$ grande devant $\pi$,} \deleted{et $f$ cofinal, alors $I'$ est grande devant $\pi$}

$(a_\pi)$ $I$ gde \unsure{dev.t|dev.} $\pi$, $f$ cofinal

$(b_\pi)$ $I'$ grande devant $\pi$, les $\vphantom{I}_{i'}\backslash I$ aussi.

Alors $a_\pi\Rightarrow b_\pi$, et si $f$ fidèle, $a_\pi$ et $b_\pi$ st équiv.
% archive sheet number: 307
